% =====================================================================
%  3. CODE  -  Người viết: TV1
%  Bản nháp: TV1 đọc lại, sửa bằng lời của mình trước khi gửi TV5.
% =====================================================================
\section{Code}\label{sec:code}\writer{TV1}

The project was implemented in Python instead of MATLAB, as recommended in the
assignment, because Python is open source and free of licensing issues. The
whole program is a single script, \texttt{lab03.py}, listed in
Appendix~\ref{app:code}. It keeps the variable names of the original MATLAB
project (\texttt{ImJPG}, \texttt{redChannel}, \texttt{ImJPG\_Gray}, \dots) so
that each result can be matched with the corresponding task.

\subsection{Environment}
The code needs Python~3 and three libraries:
\begin{itemize}
  \item \textbf{NumPy}~\cite{harris2020numpy}: arrays and matrix operations
        (\texttt{@}, \texttt{linalg.det}, \texttt{linalg.inv}, \texttt{linalg.cond});
  \item \textbf{Pillow}: reading the image file \texttt{daylilies.jpg};
  \item \textbf{Matplotlib}~\cite{hunter2007matplotlib}: displaying the images.
\end{itemize}
They are installed and the program is run with:
\begin{lstlisting}[language=bash,numbers=none]
pip install numpy matplotlib pillow
python lab03.py      # daylilies.jpg must be in the same folder
\end{lstlisting}

\subsection{Program structure}
The script has two parts. The first part imports the libraries and defines
three shared functions used by every task (\texttt{to\_uint8},
\texttt{apply\_filter\_loop}, \texttt{show}). The second part contains the
14 tasks in order; each task starts with a cell marker \texttt{\# \%\% Task k},
so it can be run on its own in VS Code or Spyder. The answers to questions
Q1--Q6 are written as comments directly under the related task. All figures
are shown at once by \texttt{plt.show()} at the end of the file.

Since the seven members wrote their tasks in separate files
(\texttt{part1.py}, \dots, \texttt{part7.py}), a small helper script,
\texttt{ghep\_file.py}, merges them into \texttt{lab03.py}: it keeps the shared
part only once, puts the tasks in the order 1 to 14 (Task 12 needs a variable
from Task 11, and Task 13 one from Task 5), and checks that every variable
required by the assignment exists.

Table~\ref{tab:matlab-python} lists the main MATLAB commands of the project
and the Python code used instead.

\begin{table}[H]\centering\small
\caption{MATLAB commands of the project and their Python equivalents.}
\label{tab:matlab-python}
\begin{tabular}{lll}
\toprule
MATLAB & Python & Remark \\
\midrule
\texttt{imread('daylilies.jpg')} & \texttt{np.array(Image.open(...).convert("RGB"))} & array $m\times n\times 3$, uint8 \\
\texttt{[m,n,l] = size(ImJPG)} & \texttt{m, n, l = ImJPG.shape} & \\
\texttt{ImJPG(:,:,1)} & \texttt{ImJPG[:, :, 0]} & indices start at 0 \\
\texttt{A*B} & \texttt{A @ B} & \texttt{A * B} is element-wise \\
\texttt{for i=1:m ... end} & \texttt{for i in range(m):} & loop over pixels \\
\texttt{reshape(x,3,1)} & \texttt{x.reshape(3, 1)} & column vector \\
\texttt{double(x)} & \texttt{x.astype(np.float64)} & \\
\texttt{uint8(x)} & \texttt{to\_uint8(x)} & own function \\
\texttt{inv, det, cond} & \texttt{np.linalg.inv, det, cond} & \\
\texttt{figure; imshow; title} & \texttt{show(img, title)} & own function \\
\bottomrule
\end{tabular}
\end{table}

\subsection{Key functions}

\paragraph{\texttt{to\_uint8(x)}.}
MATLAB's \texttt{uint8} rounds a number to the nearest integer and saturates
it to $[0,255]$, e.g.\ \texttt{uint8(260)} gives 255. In NumPy,
\texttt{.astype(np.uint8)} does not saturate but wraps around modulo 256
(260 becomes 4), which would create wrong colored spots in bright areas.
The function therefore rounds halves away from zero, as MATLAB does, and
clips the result explicitly:
\begin{lstlisting}
def to_uint8(x):
    x = np.sign(x) * np.floor(np.abs(x) + 0.5)
    return np.clip(x, 0, 255).astype(np.uint8)
\end{lstlisting}

\paragraph{\texttt{apply\_filter\_loop(img, M)}.}
A filter multiplies the color vector $p=[R;G;B]$ of every pixel by a
$3\times3$ matrix $M$. Following the project text, this is done with two
nested loops over the pixels: for each pixel the color vector is taken,
converted to floating point (MATLAB \texttt{double}), reshaped to a
$3\times1$ column (MATLAB \texttt{reshape}), multiplied by the filter matrix
and converted back with \texttt{to\_uint8} (MATLAB \texttt{uint8}):
\begin{lstlisting}
def apply_filter_loop(img, M):
    m, n, l = img.shape
    out = np.zeros((m, n, l), dtype=np.uint8)
    for i in range(m):
        for j in range(n):
            PixelColor = img[i, j, :].astype(np.float64).reshape(3, 1)
            out[i, j, :] = to_uint8(M @ PixelColor).ravel()
    return out
\end{lstlisting}
The pixel values are converted to floating point before the product, so that
no intermediate result is rounded or overflows. All matrix filters in
\texttt{lab03.py} (Tasks 4--8 and 10--13) use this single function. Filtering
the $308\times322$ image takes about one second.

\paragraph{\texttt{show(img, title)}.}
Opens a new figure and displays the image with a title. A single-channel image
(for example \texttt{redChannel}) is shown in gray levels.

\subsection{Differences from the MATLAB version and verification}
Besides the indexing from 0 and the operator \texttt{@}, two tasks needed extra
care. In Task~9 the negative is computed as \texttt{255 - ImJPG}, which never
leaves $[0,255]$. In Task~14 MATLAB adds 30 to a \texttt{uint8} image with
saturation, while NumPy would wrap around (250 + 30 = 24); the addition is
therefore done in floating point and followed by \texttt{to\_uint8}.

The results were checked against the values expected from the project:
the image size is $m=308$, $n=322$, $l=3$; $\det(\texttt{UserMatrix}) = 0.3025$;
restoring \texttt{ImJPG\_User} with the inverse filter differs from the
original image by at most 1 (out of 255); $\det(\texttt{SepiaMatrix})\approx 1.21\times10^{-7}$
and $\operatorname{cond}(\texttt{SepiaMatrix})\approx 16\,659$. The figures of
Tasks 3 and 7 match Fig.~5 and Fig.~6 of the project text.
